\section*{Statement of Significance}

Survey programmes are regularly asked about places they did not reach. The
usual route borrows strength from the places it did reach, which requires knowing
how much of the observed variation between places is real and how much is
sampling noise. That split must be estimated, often from small samples, and
practitioners have little guidance on when the resulting intervals can be
trusted.

This paper shows that the question usually asked---how precisely must we know the
variation between places?---is the wrong one. The interval depends on the
\emph{ratio} of the two variances, and a requirement stated on the variation
alone can overstate the number of places needed many times over. We also show why
two similar methods can swap places, one narrower in one survey and wider in
another, and identify what decides it.

The practical products are an interval computed rather than searched for, with a
guarantee stated under explicit conditions; an account of where it stops working;
and a diagnostic an analyst can compute from data already in hand. We are equally
explicit about what is not established, above all the gap between the model in
which the guarantee is proved and the complex designs of real surveys.
